\subsection*{3.2. Picard--Lefschetz 公式（超越情形）}

\subsubsection*{3.2.1.}
对映射
\[
f=\sum_{i=1}^{n+1}z_i^2:\mathbb C^{n+1}\longrightarrow\mathbb C.
\]
我们将把 3.1 的结果具体写出。对 $\delta$ 取函数 $\sum |z_i|^2$，并取由
\[
\left|\sum z_i^2\right|\leq a^2,
\qquad
\sum |z_i|^2\leq b^2,
\]
定义的盒形区域 $B=B_{a,b}$，其中 $0<a<b$。使用闭盒形区域而非 3.1.2 中的盒形区域，只是为了简化记号。

\subsubsection*{3.2.2.}
消灭圈流形 $V$ 的方程为
\begin{equation}\tag{3.2.2.1}
\begin{array}{rcl}
\displaystyle\sum z_i^2&=&a^2,\\
\displaystyle\sum |z_i|^2&\leq&b^2.
\end{array}
\end{equation}
将 $z=(z_i)$ 用其实部与虚部写成 $z=x+iy$。前述方程化为
\begin{equation}\tag{3.2.2.2}
\begin{array}{rcl}
x^2-y^2&=&a^2,\\
x\mathbin{\cdot}y&=&0
  \qquad\bigl(\text{即 }x\perp y\text{ 位于 }\mathbb R^{n+1}\bigr),\\
x^2+y^2&\leq&b^2.
\end{array}
\end{equation}
\footnote{原印本在 (3.2.2.2) 中写作 $\mathbb R^n$，但 $z$ 有 $n+1$ 个坐标。英文文本将这一明显的维数笔误改为 $\mathbb R^{n+1}$。}
在 $V$ 上引入新坐标系
\[
X=\frac{x}{|x|},
\qquad
Y=y\left(\frac{b^2-a^2}{2}\right)^{-1/2}.
\]
于是 $x=X(a^2+y^2)^{1/2}$，且 (3.2.2.2) 化为
\begin{equation}\tag{3.2.2.3}
\begin{array}{rcl}
X^2&=&1,\\
X\mathbin{\cdot}Y&=&0,\\
Y^2&\leq&1.
\end{array}
\end{equation}
这将 $V$ 等同于球面 $S^n$ 的切球丛。

\subsubsection*{3.2.3.}
映射
\[
\theta\longmapsto \exp(2\pi i\theta)a^2
\]
将 $\mathbb R$ 等同于 $\mathbb C$ 中以 $0$ 为中心、半径为 $a^2$ 的圆周 $S_{a^2}$ 的万有覆盖。丛 $B$ 到 $\mathbb R$ 的拉回同构于 $\mathbb R\times V$。
我们通过从 $V$ 到 $B$ 在 $\exp(2\pi i\theta)a^2$ 上的纤维的同构
\[
v\longmapsto \exp(\pi i\theta)v
\]
将其平凡化。设 $T_0$ 是此平凡化按 3.1.11 给出的映射。在坐标 (3.2.2.3) 中，
\[
T_0(X,Y)=(-X,-Y).
\]
\begin{editorialnote}
原文在 3.2.1、3.2.2 采用 $f=a^2$，但 3.2.3、3.2.4 的基点及圆盘半径写作 $a$。这里统一保留前面的 $f=a^2$ 约定，将后两节的基点、圆周及圆盘半径写为 $a^2$；乘以 $e^{\pi i\theta}$ 的平凡化公式不变。
\end{editorialnote}

\subsubsection*{3.2.4.}
作为“丛” $B$ 的边界的、半径为 $a^2$ 的圆盘 $D_{a^2}$ 上的丛，其方程为
\[
\sum |z_i|^2=b^2,
\qquad
f(z)\in D_{a^2},
\]
到圆盘的投影为 $f$。下面描述此丛在每条射线
\[
e^{2\pi i\theta}\mathbb R^+\cap D_{a^2}.
\]
上的逆像的一个平凡化。

在射线 $\mathbb R^+\cap D_{a^2}$ 上，取到 $(\mathbb R^+\cap D_{a^2})\times\partial V$ 的同构，它在 $\partial V$ 上的坐标 (3.2.2.3) 中由
\[
z=x+iy\longmapsto(f(z),X,Y)
  =\left(f(z),\frac{x}{|x|},\frac{y}{|y|}\right).
\]
给出。在 $(e^{2\pi i\theta}\mathbb R^+)\cap D_{a^2}$ 上，使用此平凡化在乘以 $e^{\pi i\theta}$ 下的像。应用 3.1.13，并在 $\partial V$ 上使用满足
\begin{equation}\tag{3.2.4.1}
X^2=Y^2=1,
\qquad
X\mathbin{\cdot}Y=0.
\end{equation}
的坐标 $(X,Y)$。若 $\varphi(\theta)$ 将 $(X,Y)$ 映至 $(X_\theta,Y_\theta)$，则算得
\begin{equation}\tag{3.2.4.2}
X_\theta+iY_\theta
=e^{\pi i\theta}(X+iY)
\qquad(0\leq\theta\leq1).
\end{equation}

\subsubsection*{3.2.5.}
总之，消灭圈流形 $V$ 是 $S^n$ 的切球丛。单值变换由 $(T_0,\varphi)$ 表示，其中 $T_0$ 是对径映射的微分，而 $\varphi$ 是从 $\operatorname{Id}$ 到 $\partial T_0$ 的如下道路：在 $\varphi(\theta)$ 下，将每个单位切向量沿它所确定的测地线平行移动距离 $-\pi\theta$（负号意为“沿向量所指的反方向”）\EditorialFootnote{原文此处写 $+\pi\theta$。但 (3.2.4.2) 给出 $X_\theta=\cos(\pi\theta)X-\sin(\pi\theta)Y$，故基点的初速度为 $-\pi Y$，沿单位切向量的反方向运动。负号也与紧接着给出的 $-\pi(1-\varphi(|u|))$ 的单值变换公式相符。}。在 $\theta=1$ 时到达对径向量。

由 3.1.12，得到 $T$ 的如下描述。\EditorialFootnote{原文这里引用 3.1.13。由有序对 $(T_0,\varphi)$ 加入领圈修正以构造 $T$ 的公式实际在 3.1.12；3.1.13 给出的是边界道路同伦类的计算变体。}设 $\varphi:[0,1]\to[0,1]$ 是 $C^\infty$ 映射，在 $0$ 附近为 $0$，在 $1$ 附近为 $1$。映射 $T$ 将 $u$ 沿 $u$ 所确定的测地线移动距离
\[
-\pi\bigl(1-\varphi(|u|)\bigr).
\]
在 $|u|=1$ 附近它为恒等映射；在 $|u|=0$ 附近它为对径映射。